\documentclass[11pt]{article}
\usepackage[margin=1in]{geometry}
\usepackage{amsmath,amssymb,amsthm,booktabs,array}
\usepackage[hidelinks]{hyperref}
\usepackage{xurl}
\setlength{\emergencystretch}{3em}
\newtheorem{theorem}{Theorem}
\newtheorem{lemma}[theorem]{Lemma}
\newtheorem{proposition}[theorem]{Proposition}
\newtheorem{corollary}[theorem]{Corollary}
\newcommand{\mr}{\operatorname{mr}}
\newcommand{\ar}{\operatorname{ar}}
\newcommand{\Z}{\mathbb Z}
\newcommand{\F}{\mathbb F}
\newcommand{\Sig}{\Sigma_2}
\title{Cyclic branched-cover homology and torsion generators\protect\\
for a prime four-bridge family}
\author{}
\date{18 November 2026}
\begin{document}\maketitle
\begin{abstract}
An explicit prime four-bridge family with complex augmentation rank
three has cyclic double-cover homology of order $117+366n$.
The calculation uses one marked filling parent and an integral
abelianization, without a sequence of sampled fillings. Every knot
exterior is hyperbolic. Its order-two orbifold group needs exactly
four torsion generators, although Coxeter reflection quotients give
bounds at most three. Consequently any link whose orbifold group
surjects onto it has at least four bridges. A current orbifold
domination theorem also excludes every arborescent knot as a source
of such an epimorphism. Cyclic homology rules out the nonunit
Lustig--Moriah plat criterion uniformly, including parameters for
which a prime cube divides the determinant. We compare only precise
published sufficient criteria and retain their unrestricted boundaries.
\end{abstract}

For $n\in\Z_{\ge0}$, let $K_n$ be the closure of the four-braid
\begin{equation}\label{eq:braid}
\beta_n=\sigma_1\sigma_2\sigma_3^{-2}\sigma_2\sigma_3^{-2}
\sigma_2\sigma_1^2\sigma_2^3\sigma_3^{-1}\sigma_2^{3+6n}.
\end{equation}
We use two established results for this exact family.
The accompanying family theorem \cite[Theorem \texttt{thm:main},
Lemma \texttt{lem:alexander}]{family} proves primeness and
\[
 \mr(K_n)=b(K_n)=4,\quad \ar(K_n,\mathbb C)=3,
 \quad g(K_n)=7+3n,\quad c(K_n)=17+6n.
\]
It supplies a full meridional $A_8$ quotient with double-transposition
meridians and excludes Coxeter reflection quotients of rank at least
four, including infinite Coxeter groups.
The accompanying cover theorem \cite[Theorems \texttt{thm:binding}
and \texttt{thm:allcovers}]{covers} identifies the actual covers as
\begin{equation}\label{eq:fill}
                \Sig(K_n)\cong N(-1,2+6n)
\end{equation}
in one fixed peripheral basis $(m,l)$, and proves that every cover is
closed hyperbolic. Here $N$ is the specified marked fourteen-tetrahedron
outside-tangle cover. Its topological identification and surviving
peripheral marking are exact; a sampled isometry does not choose them.
These family and cover results are explicit premises throughout.

The new homology calculation determines the whole integral group,
rather than just its order. This is useful because a cyclic group can
have order divisible by arbitrarily high prime powers while retaining
at most one mod-$p$ generator. A second argument retains the actual
covering involution: an invariant hyperbolic metric makes its quotient
a cone orbifold. This both gives hyperbolic knot exteriors and identifies
all orbifold torsion with the meridian class. The existing $A_8$ quotient
then gives a sharp torsion-generation bound which survives arbitrary
orbifold epimorphisms.

\section{Integral homology of the fixed parent}

\begin{lemma}\label{lem:parent}
In the marking of \eqref{eq:fill}, $H_1(N;\Z)=\Z x$ and
\[
                         [m]=5x,\qquad [l]=61x.
\]
\end{lemma}
\begin{proof}
The exact dual-cell presentation of the marked parent has generators
$a,b,c,d,e,f,g,h,i,j,k$ and the following ten relators:
\[
\begin{gathered}
 aGG,\quad eag,\quad HddA,\quad fbD,\quad ibKJ,\\
 biFH,\quad cJb,\quad kde,\quad ckhEf,\quad CJBi.
\end{gathered}
\]
A capital letter denotes the inverse of its lowercase generator.
The marked peripheral words are
\[
                       m=kDkdfbcckdfb,\qquad l=BFDKCC.
\]
These literal words are reconstructed from the same retained marked
triangulation used in the cover theorem. The pinned exact cell and
peripheral-word algorithms of SnapPy~3.3.2 are the combinatorial
import \cite{snappy}; no presentation simplification is invoked.

Here is a direct integral calculation from those words. Write the
abelianized generators additively. Unit-coefficient eliminations give
\[
\begin{gathered}
 a=2g,\quad e=-3g,\quad h=2d-2g,\quad f=d-b,\quad k=-d+3g,\\
 j=4d-b-5g,\quad c=j-b,\quad i=3d-2b-2g.
\end{gathered}
\]
The last two unused relations become
\[
                     g=6d-3b,\qquad 8g=5d.
\]
Thus the remaining presentation is
$\langle b,d\mid 43d-24b=0\rangle$.
Put $x=9d-5b$. The identities
\[
 43x-b=9(43d-24b),\qquad 24x-d=5(43d-24b)
\]
give $b=43x$ and $d=24x$. Conversely the assignments $b\mapsto43$,
$d\mapsto24$ define a homomorphism to $\Z$ sending $x$ to one.
Hence $x$ generates an infinite cyclic group; no nonunit has been
inverted in the elimination. The eleven generator images are
\[
 (30,43,-65,24,-45,-19,15,18,-44,-22,21)x.
\]
Substitution in the two literal peripheral words gives $5x$ and
$61x$, as claimed.
An exact supplement independently verifies this using mutually inverse
integral change-of-basis matrices and a unit diagonal relation block.
Its optional topology replay reconstructs all relators and peripheral
words from the marked parent, and seven corrupted controls are rejected.
The finite integer replay alone does not reprove the topological
presentation import.
\end{proof}

\begin{theorem}\label{thm:homology}
For every $n\ge0$,
\[
                H_1(\Sig(K_n);\Z)\cong\Z/(117+366n)\Z.
\]
Consequently $\dim_{\F_p}H_1(\Sig(K_n);\F_p)\le1$ for every prime $p$.
\end{theorem}
\begin{proof}
For a connected manifold with one torus boundary, filling a primitive
slope $\alpha$ gives $H_1(N(\alpha);\Z)=H_1(N;\Z)/\langle[\alpha]\rangle$.
Indeed Mayer--Vietoris presents the filling homology as
$(H_1(N)\oplus\Z[\text{core}])/\operatorname{im}H_1(T^2)$.
Choose a complementary primitive curve to $\alpha$. Its relation
eliminates the core generator, leaving exactly $[\alpha]=0$;
the $H_0$ map is injective because both pieces and their intersection
are connected. There is no extra free summand.
By Lemma~\ref{lem:parent}, the slope in \eqref{eq:fill} has image
\[
                    -5x+(2+6n)61x=(117+366n)x.
\]
This proves the integral formula uniformly. The field statement follows
by tensoring with $\F_p$; $H_0$ has no torsion term.
\end{proof}

\section{The actual knot exteriors and orbifold torsion}

Set
\[
 \Gamma_n=\pi_1(S^3\setminus K_n)/
                       \langle\!\langle\mu^2\rangle\!\rangle.
\]
This is the fundamental group of the orbifold $\mathcal O_n$ with
underlying pair $(S^3,K_n)$ and branching order two. Define
$t(G)$ to be the least number of finite-order elements generating
$G$, whenever such a generating set exists. Identity elements do
not reduce this number. In the present groups every nonidentity
torsion element will have order two, so this is also the invariant
$m_2$ discussed in \cite[Section~8]{bkn}.

\begin{lemma}\label{lem:orbifold}
For every $n\ge0$, the actual orbifold $\mathcal O_n$ is a compact
orientable hyperbolic orbifold, with singular locus $K_n$ of order two.
\end{lemma}
\begin{proof}
Let $Y_n=\Sig(K_n)$, retaining its actual smooth covering involution
$\tau_n$. It is a closed hyperbolic manifold by the cover theorem.
Dinkelbach--Leeb's Theorem~H \cite[pp.~1129, 1166--1167]{dl}
conjugates this finite smooth action to an isometric action.
The induced quotient homeomorphism preserves the branch pair, so its
quotient is the actual $(S^3,K_n)$, rather than an unspecified knot
with the same unmarked covering manifold. Since the action is
orientation preserving and has order two, its local model along the
fixed link is a half-turn. The quotient is therefore a compact
orientable hyperbolic cone manifold with cone angle $\pi$ and singular
locus $K_n$. This is a complete hyperbolic orbifold structure on
$\mathcal O_n$.
\end{proof}

\begin{theorem}\label{thm:exterior}
For every $n\ge0$, $S^3\setminus K_n$ admits a complete
finite-volume hyperbolic metric.
\end{theorem}
\begin{proof}
By Lemma~\ref{lem:orbifold}, the actual $(S^3,K_n)$ has a compact
orientable hyperbolic cone structure of angle $\pi$.
Kojima's Theorem~1.2.1 \cite[pp.~5--6]{kojima} states that the
nonsingular part of an orientable finite-volume hyperbolic cone
three-manifold has a complete finite-volume hyperbolic structure.
Apply it to this actual quotient. Its nonsingular part is
$S^3\setminus K_n$, proving the exterior assertion.
The isometric-action and cone-complement theorems
are imported at their stated scopes. This argument does not assert
that the original quotient metric restricted to the complement is
already complete.
\end{proof}

\begin{theorem}\label{thm:torsion}
Every nonidentity torsion element of $\Gamma_n$ is conjugate to the
meridian image and has order two. Moreover
\[
                              t(\Gamma_n)=m_2(K_n)=4.
\]
\end{theorem}
\begin{proof}
Realize $\mathcal O_n$ as $\mathbb H^3/\Gamma_n$ using the
isometric-action argument in Lemma~\ref{lem:orbifold}.
A finite-order isometry fixes a point of $\mathbb H^3$.
For example, in the hyperboloid model, sum the finitely many orbit
vectors of any future unit timelike vector and normalize the sum;
the future timelike cone is convex, so the normalization is defined
and fixed. All point stabilizers in this orbifold are trivial or
cyclic of order two. A nonidentity torsion element consequently
belongs to such a local order-two stabilizer. Since the singular
locus is the connected knot $K_n$, all these local generators are
conjugate to the meridian image.

The four bridge meridians generate $\Gamma_n$, giving $t\le4$.
The family theorem supplies a full $A_8$ quotient taking their images
to double transpositions. It factors through $\Gamma_n$ because
meridian squares vanish in $A_8$. By the first paragraph, every
nonidentity torsion element maps to a conjugate of the same double
transposition. If $r$ double transpositions generate a transitive
group on eight letters, their transposition-edge graph is connected:
its components are exactly the orbits of the group they generate.
It needs at least seven edges, while the $r$ generators supply at
most $2r$. Thus $r\ge4$. Any torsion generating set of $\Gamma_n$
would map onto such a generating set of $A_8$, proving $t\ge4$.
This argument needs neither a Coxeter quotient nor an ordinary
rank-four bound on $A_8$.
\end{proof}

For comparison, the independent family theorem gives
$\ar(K_n,\mathbb C)=3$ and bounds every Coxeter reflection-class
certificate by three. The new torsion-generation statement concerns
$\Gamma_n$; it does not compute its unrestricted ordinary group rank.

\section{Orbifold epimorphisms and arborescent sources}

For a link $L$, write $\Gamma(L)$ for its order-two orbifold group.
Its canonical parity map $\nu_L:\Gamma(L)\to C_2$ sends every
meridian image to the nontrivial element. We use the current
Boileau--Kitano--Nozaki definition \cite[Definition~1.1]{bkn}:
$L\succeq L'$ means there is a surjection
$\varphi:\Gamma(L)\to\Gamma(L')$ satisfying
$\nu_{L'}\varphi=\nu_L$. This parity condition is part of the
version-three definition and must not be suppressed for links.

\begin{corollary}\label{cor:bridge}
For every link $L$, any abstract group epimorphism
$\Gamma(L)\twoheadrightarrow\Gamma_n$ implies $b(L)\ge4$.
\end{corollary}
\begin{proof}
A $b(L)$-bridge presentation generates $\Gamma(L)$ by that many
meridian involutions. Their images under any homomorphism have
finite order or are the identity. If the homomorphism is surjective,
they generate $\Gamma_n$, so Theorem~\ref{thm:torsion} gives
$b(L)\ge t(\Gamma_n)=4$. Parity or meridian preservation of the
homomorphism is unnecessary here.
\end{proof}

\begin{corollary}\label{cor:arborescent}
No arborescent link $L$ with $\det L\ne0$ satisfies $L\succeq K_n$.
In particular there is no abstract orbifold-group epimorphism from
any arborescent knot to $K_n$. The first assertion also holds with
$K_n\#(\#^s3_1)$ as target, for every $s\ge0$.
\end{corollary}
\begin{proof}
Boileau--Kitano--Nozaki's Theorem~1.7 \cite{bkn} says that, for an
arborescent $L$ with nonzero determinant and a nontrivial
$\pi$-dominated target, every prime factor of the target's double
cover has a Seifert fibered JSJ piece. Our target cover is closed
hyperbolic, so it has no such piece. For the trefoil sums the lifted
sum sphere gives the prime factor $\Sig(K_n)$ in the double cover,
which yields the same contradiction.
For a knot $J$, $\Gamma(J)^{\rm ab}=C_2$ and $\nu_J$ is its unique
nontrivial map to $C_2$. An epimorphism between knot orbifold groups
induces the nontrivial map on these abelianizations, so it automatically
commutes with parity. Knot determinants are nonzero. The link
assertion therefore applies to every abstract epimorphism from an
arborescent knot.
\end{proof}

The determinant condition cannot be dropped under the broad
arborescent convention in \cite{bkn}, which includes graph prime
sums. Let $U_4$ be the four-component unlink. Its orbifold group is
$C_2*C_2*C_2*C_2$ and its double cover is
$\#^3(S^2\times S^1)$, so it is arborescent in that convention and
has determinant zero. Sending the four free involutions to the
four bridge meridians of $K_n$ gives a parity-preserving epimorphism
$\Gamma(U_4)\twoheadrightarrow\Gamma_n$. This also realizes equality
in the bridge lower bound of Corollary~\ref{cor:bridge}.

\section{Uniform comparison with specified equality criteria}

\begin{proposition}\label{prop:elementary}
For all $n\ge0$, $K_n$ is nonalternating and
$s(K_n)=6+6n$. Neither chirality is positive or strongly quasipositive.
No $K_n$ is a genuine cable on a nontrivial companion, a braid-pattern satellite of index at least two on a nontrivial companion, or a twisted or untwisted Whitehead double.
\end{proposition}
\begin{proof}
The family Alexander formula is
\[
 \Delta_n(t)=\frac{P(t)+t^{6n+6}t^9P(t^{-1})}{1+t},\quad
 P(t)=1-3t+4t^2-8t^4+18t^5-21t^6+15t^7-6t^8+t^9.
\]
The second numerator term starts in degree at least six.
Thus the coefficients through degree five are always
$1,-4,8,-8,0,18$. The zero is between nonzero extreme coefficients,
contradicting Murasugi's nonvanishing theorem for alternating knots
\cite[Theorem~1.1]{murasugi}.
The displayed braid is homogeneous, has writhe $7+6n$, four Seifert
circles and two components after smoothing all negative crossings.
Abe's formula \cite[Theorem~1.2]{abe} gives
$s=7+6n-4+4-1=6+6n$; the mirror has $s=-6-6n$.
Since $2g=14+6n$, neither chirality has $s=2g$.
Abe's Theorem~1.3 excludes positive and strongly quasipositive
presentations for these homogeneous knots. Arbitrary quasipositivity
is not excluded by this statement.

Hyperbolic exteriors already exclude genuine satellites.
There are also independent elementary checks for the displayed
classes. A winding-$p$ cable or $p$-strand braid satellite on a
nontrivial companion has bridge number at least $p\,b(J)$
\cite[Theorem~1]{schultens}. Thus $p=2$; a two-strand braid pattern
is a $(2,q)$ cable. The satellite Alexander formula
\cite[Theorem~2 and Seifert corollary]{lm} gives
$|q|=\det K_n=117+366n$, since the companion factor is evaluated
at one. The torus-pattern factor alone has breadth
$|q|-1=116+366n>14+6n=2g(K_n)$, impossible because Laurent
polynomial breadths add under multiplication. A Whitehead pattern,
with any twisting, bounds a once-punctured torus: take the companion
annulus and attach its clasp band. Its genus is at most one,
whereas $g(K_n)\ge7$. Iterated doubles are covered by their outermost
Whitehead pattern. Primeness separately excludes nontrivial
connected-sum constructions.
\end{proof}

\begin{lemma}\label{lem:fox}
For any knot diagram $D$ of $K$ and prime $p$,
\[
 \dim_{\F_p}\operatorname{Col}_p(D)
       =1+\dim_{\F_p}H_1(\Sig(K);\F_p).
\]
\end{lemma}
\begin{proof}
We recall the direct group calculation behind the classical core-group
interpretation \cite{core}. In the order-two orbifold group put
$a=\mu_0$ and rewrite the parity kernel with transversal $\{1,a\}$.
The square relators pair the Schreier generators
$u_i=\mu_i a^{-1}$ and $v_i=a\mu_i$ by $u_iv_i=v_iu_i=1$;
$u_0=1$ and the base square gives $v_0=a^2=1$.
Eliminate $v_i=u_i^{-1}$ and write $x_i=u_i=\mu_i a$, $x_0=1$.
The crossing relator
$\mu_{\rm out}\mu_{\rm over}\mu_{\rm in}\mu_{\rm over}$ has rewrites
\[
 R=x_{\rm out}x_{\rm over}^{-1}x_{\rm in}x_{\rm over}^{-1},
 \qquad x_{\rm out}^{-1}R^{-1}x_{\rm out},
\]
so only the first is needed. These exhaust the rewritten relators;
no relation $x_i^2=1$ is introduced. The parity kernel is
$\pi_1(\Sig(K))$, so its abelianization is presented by $x_0=0$
and $2x_{\rm over}-x_{\rm in}-x_{\rm out}=0$.
Its homomorphisms to $\F_p$ are precisely Fox colourings with base
colour zero. Subtracting the constant base colour splits off one
copy of $\F_p$. The remaining dual homology space has dimension
$\dim H_1(\Sig(K);\F_p)$, proving the formula, including $p=2$.
\end{proof}

\begin{theorem}\label{thm:criteria}
No $K_n$, for any $n\ge0$, satisfies either of the following
sharp-four sufficient criteria:
\begin{enumerate}
\item a construction supplying a rank-four Coxeter quotient with
meridians sent to reflections;
\item the nonunit Lustig--Moriah matrix-gcd condition in any plat presentation,
or its geometric generalized-Montesinos criterion with nonunit
numerator gcd.
\end{enumerate}
It also cannot be an output whose equality certificate requires
ordinary complex augmentation rank four.
\end{theorem}
\begin{proof}
The first and last assertions follow from the separate uniform
Coxeter and augmentation bounds. For the second, let $R$ be the
unreduced Burau matrix at $t=-1$ for a $2b$-plat. In the conventions
of \cite[Proposition~5.6, pp.~186--187]{lustig}, also restated in
\cite[Theorem~4.20]{moriah}, let the $i$th row of $L$ be
$e_{2i-1}^{T}-e_{2i}^{T}$ and the $j$th column of $M$ be
$e_{2j-1}+e_{2j}$. Suppose the entries of $LRM$ have nonunit gcd.
A prime $p$ divides all entries (choose any prime if all are zero).
Every $2b$-plat gives a bridge upper bound $b$, so any such
presentation of $K_n$ has $b\ge4$.

For completeness, the Fox-colouring consequence gives the needed
obstruction directly. The two colours in each top cap coincide,
so top colour vectors form $\operatorname{im}M$. The matrix $R$
propagates them through the braid by the Fox crossing rule, and
$L$ tests equality at the bottom cap pairs. The constraints are
$LRM v=0$.
If $LRM=0$ over $\F_p$, all $b$ cap colours are free, so the
Fox-colouring space has dimension $b$. Lemma~\ref{lem:fox} would therefore give
$\dim_{\F_p}H_1(\Sig(K_n);\F_p)=b-1\ge3$. This contradicts Theorem~\ref{thm:homology}.
The geometric nonunit numerator-gcd criterion is a specialization
of this matrix condition in the original construction; no claim
is made for geometric descriptions with unit gcd.
\end{proof}

For example $n=3$ has determinant $1215=3^5\cdot5$.
A prime-cube divisibility test alone cannot exclude it; the integral
homology is nevertheless cyclic, and the mod-three dimension is one.
The argument excludes the stated criterion in every possible
presentation satisfying it, rather than only a displayed plat.

The Coxeter-producing class in Theorem~\ref{thm:criteria} includes the
precise twisted-link and even-weight bipartite arborescent theorems
of \cite{bbk}, the many-twig arborescent theorem of \cite{twigs},
the nonunit even-column fishnets and further maximal-rank Coxeter
constructions of \cite{fishnet}, and any successful maximal-rank
finite-Coxeter search of \cite{bkm}. The integer-box classical families of
\cite{jp} also give such quotients by the written boundary-label
transport in the accompanying single-knot comparison \cite{comparison}:
replace the four step-cycle labels by the four simple reflections
with adjacent exponent $2p-1$ and nonadjacent exponent two;
the permitted two-strand twist multiples preserve the labels.
This includes rational tangles assembled from those permitted boxes;
it does not exclude arbitrary rational replacement without its
noncollapsed sharp quotient. Likewise, the augmentation statement
excludes maximal-augmentation satellite and symmetric rank-raising
outputs at bridge number four \cite{ch,cornwell}.

Every link can have a loose fishnet presentation; our obstruction is
to the stated sharp-quotient criteria. Adding an unknot to $K_n$ is
still allowed by \cite[Theorem~1.1]{fishnet}, which proves an equality
for the larger link. The plain-sphere theorem \cite[Theorem~1.1]{plain}
is a theorem for every link using restricted geometric presentations;
it supplies no unrestricted meridional-rank lower bound that would
be contradicted here. No exclusion of those universal statements
is asserted. Nor is unit-gcd generalized-Montesinos membership
classified by the argument above. The general matrix/Fox-ideal rank
criterion of \cite[Theorem3.5]{moriah} is also separate from the
scalar nonunit Burau condition. The class generated from arbitrary
meridionally tame base knots in \cite[Theorem1]{bdjw} includes its
base knots themselves. Hyperbolicity excludes its genuine satellite
and connected-sum outputs, but it does not decide whether a $K_n$
is meridionally tame; bridge equality alone does not imply tameness.

Strict augmentation gaps already occur for three-bridge prime knots
in \cite[Corollary~1.12]{cornwell}. The Whitehead double of the square
knot gives another satellite example with augmentation rank
at most five and meridional rank six \cite[Theorem~1.3]{whitehead}.
Thus no first-prime-gap claim is made. The distinctions established
here are the explicit four-bridge family, cyclic closed-hyperbolic
cover homology, exact orbifold torsion-generation rank, and the
specified uniform construction and domination obstructions.
A bounded source comparison does not establish worldwide priority.

\paragraph{AI assistance and review.}
ChatGPT (OpenAI) assisted with mathematical arguments, literature checks,
exposition, proof review and verification code.

\begin{thebibliography}{99}\small\raggedright
\bibitem{family} \emph{Prime four-bridge knots with augmentation rank
three}, accompanying family theorem, Theorem \texttt{thm:main},
Lemma \texttt{lem:alexander} and Corollary \texttt{cor:coxeter}.
\bibitem{covers} \emph{Hyperbolic double covers of a prime four-bridge
twist family}, accompanying marked-cover theorem,
Theorems \texttt{thm:binding} and \texttt{thm:allcovers}.
\bibitem{snappy} M. Culler, N. M. Dunfield, M. Goerner and J. R. Weeks,
SnapPy~3.3.2; exact dual-cell presentation and marked peripheral words.
\url{https://snappy.computop.org/}.
\bibitem{dl} J. Dinkelbach and B. Leeb, Equivariant Ricci flow with
surgery and applications to finite group actions on geometric
3-manifolds, Geom. Topol. 13 (2009), 1129--1173,
Theorem H, proof pp1166--1167.
\url{https://doi.org/10.2140/gt.2009.13.1129}.
\bibitem{kojima} S. Kojima, Deformations of hyperbolic 3-cone-manifolds,
J. Differential Geom. 49 (1998), 469--516.
Retained arXiv:math/9809034v1, Theorem1.2.1 and proofpp5--6.
\url{https://arxiv.org/abs/math/9809034v1}.
\bibitem{bkn} M. Boileau, T. Kitano and Y. Nozaki,
A preorder on the set of links with applications to symmetric unions,
arXiv:2510.12372v3 (19 April 2026), Definition1.1,
Theorem1.7 and Section8. \url{https://arxiv.org/abs/2510.12372v3}.
\bibitem{murasugi} K. Murasugi, On the Alexander polynomial of the
alternating knot, Osaka Math. J. 10 (1958), 181--189, Theorem1.1.
\bibitem{abe} T. Abe, The Rasmussen invariant of a homogeneous knot,
Proc. AMS 139 (2011), 2647--2656; retained arXiv:1003.5392v1,
Theorems1.2--1.3. \url{https://arxiv.org/abs/1003.5392v1}.
\bibitem{schultens} J. Schultens, Additivity of bridge numbers of knots,
Math. Proc. Cambridge Philos. Soc. 135 (2003), 539--544, Theorem1.
\bibitem{lm} C. Livingston and P. Melvin, Abelian invariants of satellite
knots, Geometry and Topology, LNM1167 (1985), 217--227, Theorem2
and its Seifert corollaryp222.
\bibitem{lustig} M. Lustig and Y. Moriah, Generalized Montesinos knots,
tunnels and $N$-torsion, Math. Ann. 295 (1993), 167--189,
Theorem0.1, Definition1.2 and Proposition5.6.
\bibitem{moriah} Y. Moriah, Heegaard splittings of knot exteriors,
Geom. Topol. Monographs12 (2007), 191--232, Theorem4.20.
\bibitem{core} D. Silver, L. Traldi and S. Williams, Core groups,
J. Knot Theory Ramifications33 (2024), 2450013, introduction
and Section7, crediting Kelly and Wada.
\bibitem{bbk} S. Baader, R. Blair and A. Kjuchukova,
Coxeter groups and meridional rank of links, Math. Ann.379 (2021),
1533--1551, Theorems1--2.
\bibitem{twigs} S. Baader, R. Blair, A. Kjuchukova and F. Misev,
The bridge number of arborescent links with many twigs,
Algebraic \& Geometric Topology23 (2023), 75--85, Theorem2.
\bibitem{fishnet} R. Blair, A. Kjuchukova and E. Pfaff, Adding a
suitable unknot to any link equates bridge number and meridional rank,
Math. Proc. Cambridge Philos. Soc.181 (2026), 1339--1355,
Proposition1.4 and Section2.2.
\url{https://doi.org/10.1017/S030500412610214X}.
\bibitem{plain} R. Blair, A. Kjuchukova and E. Pfaff,
The plain sphere number of a link, Pacific J. Math.342(2) (2026),
217--233. \url{https://doi.org/10.2140/pjm.2026.342.217}.
\bibitem{bkm} R. Blair, A. Kjuchukova and N. Morrison,
Coxeter quotients of knot groups through 16 crossings,
Experimental Math.34 (2025), 27--37.
\bibitem{jp} J. Joseph and P. Pongtanapaisan, Meridional rank and
bridge number of knotted 2-spheres, Canad. J. Math.77 (2025),282--299;
retained arXiv:2111.09233v2, Construction3.5, Definition3.6,
Corollary3.12 and Remark3.11.
\bibitem{comparison} \emph{A bounded family comparison for a
seventeen-crossing knot}, accompanying boundary-label construction,
Proposition \texttt{prop:coxeterfamilies}.
\bibitem{bdjw} M. Boileau, E. Dutra, Y. Jang and R. Weidmann,
Meridional rank of knots whose exterior is a graph manifold,
arXiv:1608.01570v1 (2016), Theorem1 and Corollary1.
\url{https://arxiv.org/abs/1608.01570}.
\bibitem{ch} C. R. Cornwell and D. R. Hemminger,
Augmentation rank of satellites with braid pattern,
Commun. Anal. Geom.24 (2016), 939--967.
\bibitem{cornwell} C. R. Cornwell, KCH representations, augmentations,
and $A$-polynomials, J. Symplectic Geom.15 (2017), 983--1017.
Retained arXiv:1310.7526v3.
\bibitem{whitehead} H. Gao, S. Gu, Z. Li and J. Wang,
Meridional ranks of iterated Whitehead doubles,
arXiv:2611.15673v1 (2026), Theorem1.3.
\end{thebibliography}
\end{document}
